\chapter{Introducción}
\label{ch:intro}

El crecimiento acelerado del IoT ha posicionado a los
sistemas embebidos de ultrabaja potencia como una pieza central de la
infraestructura tecnológica moderna. En 2025 se estimó que el número de
dispositivos IoT conectados superaba los
21\,000 millones a nivel global~\cite{iotanalytics2025},
la mayoría de los cuales operan alimentados por baterías de tamaño reducido
o mediante técnicas de recolección de energía del entorno~\cite{iotlowpowersurvey}. En este escenario,
el presupuesto de potencia de un nodo típico se sitúa en el orden de los
milivatios: el procesador, la radio, los sensores y la memoria compiten por
ese margen mínimo, y la falta de eficiencia en el firmware se traduce
directamente en una reducción de la vida útil del dispositivo.

Para responder a esta demanda, la comunidad de arquitecturas de procesadores ha
convergido en torno a RISC-V, un conjunto de instrucciones abierto, modular y
libre de regalías que permite diseñar núcleos altamente especializados para
aplicaciones de bajo consumo. Dentro de este ecosistema, el grupo de investigación PULP de ETH~Zürich ha desarrollado la plataforma PULPissimo~\cite{pulpissimo}, un
SoC orientado a sistemas embebidos que integra el procesador CV32E40P~\cite{cv32e40p} como su componente central de cómputo.

El CV32E40P implementa el ISA RISC-V RV32IMFCXpulp de 32 bits e incorpora
extensiones vectoriales SIMD y de punto flotante que lo hacen especialmente
adecuado para cargas de procesamiento de señales, control y aprendizaje
automático en el borde de la red~\cite{cv32e40p}. A pesar de sus capacidades, el núcleo carece de mecanismos de inspección energética: no expone contadores de potencia,
no implementa una interfaz interna de medición energética y no ofrece registros
dedicados que permitan obtener un perfil de ejecución útil para estimar el
consumo energético~\cite{leonvega2024efimon}.
Esta ausencia obliga a recurrir a instrumentación física externa para relacionar
una ejecución concreta con su consumo energético. Dicha instrumentación es
necesaria para caracterizar y validar modelos de consumo, pero su uso como único
mecanismo de análisis dificulta la comparación repetida de regiones de código
durante el desarrollo de firmware.

El modelado de consumo energético a nivel de instrucción es una técnica
consolidada que puede cerrar esa brecha. Propuesta originalmente por
Tiwari~\emph{et al.}~\cite{tiwari1994power} para procesadores Intel~486 y
SPARC, se basa en la observación de que cada tipo de instrucción activa un
subconjunto distinto de las unidades funcionales del procesador, generando
un perfil de consumo característico y medible. Extendida recientemente a la
arquitectura RISC-V por Fang~\cite{fang2021instruction}, esta metodología permite estimar la potencia de un programa
como la suma ponderada del número de instrucciones ejecutadas por categoría.
Sin embargo, la implementación clásica de esta metodología en la literatura
opera de forma \emph{off-chip} y \emph{post-mortem}: el análisis se realiza
sobre el binario ya generado, fuera del procesador.

Este trabajo propone integrar ese modelo en el hardware del CV32E40P mediante
contadores accesibles por CSR. Durante la ejecución bare-metal sobre
PULPissimo, el firmware ejecuta la carga de prueba y el procesador acumula el
número de instrucciones retiradas por categoría. Al finalizar la ejecución o una
ventana de medición, los registros del clasificador se leen desde la computadora
anfitriona y se combinan con coeficientes caracterizados previamente mediante un
circuito de medición de potencia. La instrumentación
externa no se elimina del flujo experimental: se usa para obtener la referencia
física y validar el modelo, mientras que los contadores internos permiten
repetir el análisis sin ejecutar una medición eléctrica completa para cada
escenario posterior.

% -----------------------------------------------------------------------
\section{Planteamiento del problema}
\label{sec:planteamiento}

\subsection{Estado actual}
\label{sec:estado_actual}

La estimación de consumo energético depende de información que el procesador no
siempre expone de forma directa. En plataformas con sistema operativo, parte de
esa información puede obtenerse mediante infraestructura de monitoreo y
herramientas externas. En una ejecución bare-metal, como la de PULPissimo sobre
el CV32E40P, esa capa no existe: el firmware se ejecuta directamente sobre el
hardware y el procesador no entrega por sí mismo conteos de instrucciones
retiradas por categoría.

Por ello, el análisis energético de esta plataforma depende de mediciones
eléctricas externas o de análisis posterior del programa ejecutado. Ambos
enfoques son útiles, pero no generan desde el propio procesador los datos
necesarios para asociar ejecución y consumo estimado dentro del flujo de
prueba.

\subsection{Definición del problema}
\label{sec:definicion}

El problema que aborda este trabajo puede enunciarse de forma precisa: el
procesador CV32E40P, ejecutando firmware en bare-metal sobre el SoC PULPissimo
sintetizado en una FPGA, no dispone de un mecanismo interno que genere un
histograma completo de instrucciones retiradas por categoría ni que permita
extraer esa información desde el flujo de prueba para estimar el consumo
energético observado en la plataforma.

Como consecuencia, un desarrollador de firmware que necesite optimizar el
consumo de su programa debe apoyarse en instrumentación eléctrica externa para
cada escenario que desea estudiar. Esta instrumentación es útil como referencia
física, pero no entrega una señal interna que pueda leerse desde el flujo de
prueba para comparar regiones de código o automatizar análisis posteriores.

Por tanto, la brecha central no es únicamente la ausencia de una medición física
de potencia, sino la falta de un mecanismo interno que relacione la ejecución
del firmware con una estimación de consumo obtenible a partir de las
instrucciones vistas durante la ejecución.

% -----------------------------------------------------------------------
\section{Estructura del documento}
\label{sec:estructura}

En el \textbf{capítulo~\ref{ch:marco}} se presenta el marco teórico, que incluye
la plataforma hardware, los fundamentos de análisis energético, el diseño RTL,
la medición eléctrica y el estado del arte relacionado con la estimación de
consumo a nivel de instrucción.

En el \textbf{capítulo~\ref{ch:solucion}} se describe la solución propuesta:
el circuito de medición de potencia, el módulo clasificador de instrucciones,
su integración en el RTL del CV32E40P, la carga y lectura mediante JTAG, y el
flujo de caracterización y análisis de datos.

En el \textbf{capítulo~\ref{ch:obj}} se enuncian la hipótesis de trabajo, los
objetivos del proyecto, las tareas asociadas, los alcances y las limitaciones.

En el \textbf{capítulo~\ref{ch:experimento}} se presenta el experimento
preliminar propuesto para validar el flujo de medición y la observabilidad de
diferencias de potencia entre rutinas dominadas por distintas categorías de
instrucción.

En el \textbf{capítulo~\ref{ch:cron}} se muestra el cronograma de ejecución
del proyecto mediante un diagrama de Gantt.
